\subsection{Уравнение второго рода со сжимающим оператором}

Рассматривается то же уравнение

\stepcounter{equation}
\begin{equation}\tag{1}\label{eq:contr-main}
  u(x) - \int\limits_{0}^{\ell} K(x,y)\, u(y)\, dy = f(x) \in \Lp{2}(0,\ell) ,
  \qquad\text{или}\qquad
  u - Au = f \ \text{ в } H ,
\end{equation}
но теперь вместо вырожденности ядра накладывается условие на норму оператора.

\begin{definition}
Оператор $A$, определённый на всём $H$, называется \textbf{сжимающим}, если
\[
  \norm{A} < 1 .
\]
\end{definition}

\begin{warn}
В тетради написано $\norm{A} \le 1$, но во всех дальнейших выкладках стоит
деление на $1 - \norm{A}$, да и сам метод последовательных приближений при
$\norm{A} = 1$ не работает. Значит неравенство строгое.
\end{warn}

\subsubsection*{Построение решения}

Метод последовательных приближений:
\[
  u_k = A u_{k-1} + f , \qquad u_0 = f ,
\]
откуда
\[
  u_m = \sum_{k=0}^{m} A^k f .
\]
Проверим фундаментальность. При $l < m$
\[
  \norm{u_m - u_l}
  = \Bigg\|\sum_{k=l+1}^{m} A^k f\Bigg\|
  \le \sum_{k=l+1}^{m} \norm{A^k}\,\norm{f}
  \le \norm{f} \sum_{k=l+1}^{\infty} \norm{A}^k
  = \norm{f}\,\frac{\norm{A}^{\,l+1}}{1 - \norm{A}}
  \xrightarrow[l,m \to \infty]{} 0 ,
\]
то есть $\{u_m\}$ фундаментальна. Пространство $H$ полно, значит
$u_m \to u \in H$. Переходя к пределу в $u_m = A u_{m-1} + f$ (оператор $A$
непрерывен, так как ограничен), получаем $u = Au + f$, то есть $u$ —
решение~\eqref{eq:contr-main}.

\subsubsection*{Единственность}

Пусть $u$ и $v$ — два решения~\eqref{eq:contr-main}, и $w = u - v$. Вычитая,
получаем $w = Aw$, а значит и $w = A^k w$ при любом $k$. Тогда
\[
  \norm{w} = \norm{A^k w} \le \norm{A}^k \norm{w} \xrightarrow[k\to\infty]{} 0 ,
\]
откуда $\norm{w} = 0$, то есть $u = v$.

\begin{idea}
Можно и в один шаг: $\norm{w} = \norm{Aw} \le \norm{A}\norm{w} < \norm{w}$
при $\norm{w} \ne 0$ — противоречие. Именно здесь и нужна строгость
неравенства $\norm{A} < 1$.
\end{idea}

\subsubsection*{Оценка обратного оператора}

Переходя к пределу в $u_m = \sum_{k=0}^m A^k f$:
\[
  \norm{(I-A)^{-1} f}
  = \Bigg\|\sum_{k=0}^{\infty} A^k f\Bigg\|
  \le \sum_{k=0}^{\infty} \norm{A}^k \norm{f}
  = \frac{\norm{f}}{1 - \norm{A}} .
\]

\begin{proposition}\label{prop:contraction}
Пусть $A$ — сжимающий оператор, действующий на всём пространстве $H$. Тогда
уравнение~\eqref{eq:contr-main} разрешимо единственным образом при любой
правой части $f \in H$, причём
\[
  \norm{(I - A)^{-1}} \le \frac{1}{1 - \norm{A}} .
\]
\end{proposition}

\begin{corollary}
Если $A$ — сжимающий оператор, то задача решения интегрального
уравнения~\eqref{eq:contr-main} \textbf{корректна}: решение существует,
единственно и непрерывно зависит от правой части, поскольку
\[
  \norm{u} \le \frac{\norm{f}}{1 - \norm{A}} .
\]
\end{corollary}

\begin{corollary}\label{cor:adj-inverse}
Если $A$ — сжимающий оператор, то
\[
  (I - A^*)^{-1} = \big[(I - A)^{-1}\big]^{*} .
\]
\end{corollary}

\begin{proof}
Так как $\norm{A^*} = \norm{A} < 1$, оператор $A^*$ тоже сжимающий, и
$(I - A^*)^{-1}$ существует. Возьмём произвольные $f, g \in H$ и положим
\[
  u = (I-A)^{-1} f \ \Longleftrightarrow\ (I-A)u = f ,
  \qquad
  v = (I-A^*)^{-1} g \ \Longleftrightarrow\ (I-A^*)v = g .
\]
Тогда
\[
  \big((I-A)^{-1} f,\ g\big)
  = (u, g)
  = \big(u,\ (I-A^*)v\big)
  = \big((I-A)u,\ v\big)
  = (f, v)
  = \big(f,\ (I-A^*)^{-1} g\big) .
\]
Равенство выполнено при всех $f, g$, а это и есть определение сопряжённого
оператора.
\end{proof}

\begin{remark}
Формально
\[
  (I - A)^{-1} = \sum_{k=0}^{\infty} A^k
\]
— это \emph{ряд Неймана}, и выкладка выше как раз и доказывает его сходимость
по операторной норме при $\norm{A} < 1$.
\end{remark}

\begin{idea}
Сравните с геометрической прогрессией $\frac{1}{1-q} = \sum_k q^k$ при
$\abs{q} < 1$: вся конструкция — её дословный операторный аналог, и оценка
$\norm{(I-A)^{-1}} \le \frac{1}{1-\norm{A}}$ тоже.

\medskip
Практический смысл: условие $\norm{A} < 1$ — это условие на «малость» ядра,
\[
  \norm{A} \le \left(\iint\limits_{D\times D} \abs{K(x,y)}^2 dx\,dy\right)^{1/2} ,
\]
и при малом ядре уравнение решается простой итерацией без всякой теории. Вся
трудность — в случае, когда ядро не мало; именно поэтому в \S6 компактный
оператор разбивается на конечномерную часть и сжимающий остаток.
\end{idea}
